\documentclass[11pt,a4paper]{article}
\usepackage[margin=1in]{geometry}
\usepackage{amsmath,amssymb,amsthm}
\usepackage{algorithm}
\usepackage{algpseudocode}
\usepackage{booktabs}
\usepackage{hyperref}

\theoremstyle{definition}
\newtheorem{definition}{Definition}[section]
\newtheorem{theorem}{Theorem}[section]
\newtheorem{lemma}[theorem]{Lemma}
\newtheorem{remark}{Remark}[section]

\title{\textbf{Rigorous Mathematical Specification, Exponential Moving Average Dynamics, and Convergence Analysis of the RMSProp Optimizer}}
\author{Team Scaibu \\ \small Email: \texttt{jaydeep.9664920749@gmail.com} \\ \small Architecture and Core Engine Team}
\date{October 2026}

\begin{document}
\maketitle

\begin{abstract}
This document provides a comprehensive theoretical and algorithmic specification of Root Mean Square Propagation (RMSProp). Designed to resolve AdaGrad's premature convergence and step-size decay in non-convex neural optimization, RMSProp replaces monotonic second-moment accumulation with an Exponential Moving Average (EMA). We formalize coordinate-wise variance tracking, present the complete pseudo-code algorithm, provide formal convergence guarantees under non-convex $L$-smooth landscapes, derive computational complexities, work through a step-by-step numerical example matching the exact reference implementation, and analyze floating-point numerical considerations.
\end{abstract}

\section{Notation and Mathematical Preliminaries}

Let $\theta \in \mathbb{R}^P$ denote the trainable parameter vector of a differentiable neural objective function $\mathcal{L}: \mathbb{R}^P \to \mathbb{R}$. Optimization proceeds across discrete iterations $t \in \{1, 2, \dots, T\}$.

\begin{itemize}
    \item $g_t \triangleq \nabla_\theta \mathcal{L}(\theta_t) \in \mathbb{R}^P$: Stochastic gradient vector evaluated at iterate $\theta_t$.
    \item $g_{t, i} \in \mathbb{R}$: The $i$-th component of the gradient vector $g_t$, for $i \in \{1, \dots, P\}$.
    \item $v_t \in \mathbb{R}^P$: Exponential moving average (EMA) of squared gradient coordinates at step $t$.
    \item $\alpha \in [0, 1)$: Exponential decay rate (smoothing factor) controlling effective memory horizon.
    \item $\eta > 0$: Base learning rate (step size).
    \item $\epsilon > 0$: Smoothing denominator constant ensuring strict numerical regularity.
    \item $\odot$ and $\oslash$: Component-wise (Hadamard) product and division operators respectively.
\end{itemize}

\subsection{Coordinate-Wise Mechanics and Memory Horizon}
The effective historical horizon $H_{\text{eff}}$ over which past squared gradients contribute is governed by the decay parameter $\alpha$:
\begin{equation}
H_{\text{eff}} = \frac{1}{1 - \alpha}
\end{equation}
For standard default $\alpha = 0.99$, $H_{\text{eff}} \approx 100$ gradient steps.

\begin{table}[h]
\centering
\caption{Summary of Mathematical Symbols for RMSProp}
\begin{tabular}{lll}
\toprule
\textbf{Symbol} & \textbf{Type / Domain} & \textbf{Description} \\
\midrule
$\theta_t$ & $\mathbb{R}^P$ & Parameter vector at time step $t$ \\
$g_t$ & $\mathbb{R}^P$ & Gradient $\nabla_\theta \mathcal{L}(\theta_t)$ \\
$v_t$ & $\mathbb{R}^P_{\ge 0}$ & Uncentered second moment vector (EMA of $g^2$) \\
$\alpha$ & $[0, 1)$ & Smoothing factor (standard default $0.99$) \\
$\eta$ & $\mathbb{R}_{> 0}$ & Learning rate (step size) \\
$\epsilon$ & $\mathbb{R}_{> 0}$ & Numerical stability constant ($10^{-8}$) \\
\bottomrule
\end{tabular}
\end{table}

\section{Problem Definition and Core Concepts}

\subsection{Intuition and Motivation}
In deep learning landscapes exhibiting ill-conditioned curvature (e.g., canyons or ravines where gradient magnitude is steep along high-curvature directions and shallow along low-curvature trajectories), standard stochastic gradient descent oscillates violently across the ravine while progressing slowly along the floor. AdaGrad stabilizes this by scaling coordinate step sizes inversely proportional to $\sqrt{\sum_{\tau=1}^t g_{\tau, i}^2}$. However, because the historical sum grows monotonically without bound, AdaGrad's effective learning rate decays strictly to zero, prematurely halting training in deep non-convex architectures. RMSProp remedies this by exponentially discounting older gradient magnitudes, allowing the effective learning rate to dynamically adapt up and down.

\subsection{Formal Mathematical Definition}
\begin{definition}[RMSProp Update Rule]
Given initial parameter vector $\theta_0 \in \mathbb{R}^P$, initial second moment buffer $v_0 = \mathbf{0} \in \mathbb{R}^P$, base step size $\eta > 0$, smoothing factor $\alpha \in [0, 1)$, and regularization constant $\epsilon > 0$, the recurrence relation for iteration $t \ge 1$ across each coordinate $i \in \{1, \dots, P\}$ is:
\begin{align}
v_{t, i} &= \alpha v_{t-1, i} + (1 - \alpha) g_{t, i}^2 \label{eq:rmsprop_v} \\
\theta_{t, i} &= \theta_{t-1, i} - \frac{\eta}{\sqrt{v_{t, i}} + \epsilon} g_{t, i} \label{eq:rmsprop_theta}
\end{align}
In vectorized vector notation:
\begin{equation}
v_t = \alpha v_{t-1} + (1 - \alpha)(g_t \odot g_t), \quad \theta_t = \theta_{t-1} - \eta g_t \oslash (\sqrt{v_t} + \epsilon)
\end{equation}
\end{definition}

\section{Algorithmic Specification}

The operational execution of RMSProp strictly mirrors the contract defined in \texttt{NnAlgoRmspropOptimizer}.

\begin{algorithm}[H]
\caption{Root Mean Square Propagation (RMSProp) Step Execution}
\begin{algorithmic}[1]
\Require Parameter vector $\theta \in \mathbb{R}^P$, gradient vector $g \in \mathbb{R}^P$, second moment buffer $v \in \mathbb{R}^P$
\Require Learning rate $\eta > 0$, decay factor $\alpha \in [0, 1)$, stability epsilon $\epsilon > 0$
\Ensure Updated parameters $\theta_{\text{next}} \in \mathbb{R}^P$, updated second moment buffer $v_{\text{next}} \in \mathbb{R}^P$
\State $P \gets \text{length}(\theta)$
\If{$P = 0$ \textbf{or} $\text{length}(g) \ne P$ \textbf{or} $\text{length}(v) \ne P$}
    \State \textbf{throw} PreconditionFailureException(``Buffer length mismatch or zero length'')
\EndIf
\If{$\eta \le 0$ \textbf{or} $\epsilon \le 0$ \textbf{or} $\alpha < 0$ \textbf{or} $\alpha \ge 1.0$}
    \State \textbf{throw} PreconditionFailureException(``Hyperparameters out of domain'')
\EndIf
\State $\theta_{\text{next}} \gets \text{new Array of size } P$
\State $v_{\text{next}} \gets \text{new Array of size } P$
\For{$i \gets 1 \textbf{ to } P$}
    \State $v_{\text{next}}[i] \gets \alpha \cdot v[i] + (1.0 - \alpha) \cdot (g[i])^2$
    \State $\theta_{\text{next}}[i] \gets \theta[i] - \left(\frac{\eta}{\sqrt{v_{\text{next}}[i]} + \epsilon}\right) \cdot g[i]$
\EndFor
\State \Return $\{\text{updated\_parameters}: \theta_{\text{next}}, \text{updated\_moving\_average}: v_{\text{next}}\}$
\end{algorithmic}
\end{algorithm}

\section{Theoretical Guarantees and Convergence Analysis}

\begin{theorem}[Convergence of Non-Convex RMSProp under Bounded Variance]
Let the objective $\mathcal{L}: \mathbb{R}^P \to \mathbb{R}$ be $L$-smooth such that $\|\nabla \mathcal{L}(\theta_1) - \nabla \mathcal{L}(\theta_2)\|_2 \le L \|\theta_1 - \theta_2\|_2$ for all $\theta_1, \theta_2$. Assume stochastic gradients are unbiased $\mathbb{E}[g_t \mid \theta_{t-1}] = \nabla \mathcal{L}(\theta_{t-1})$ with bounded coordinate variance $\mathbb{E}[(g_{t, i} - \nabla_i \mathcal{L}(\theta_{t-1}))^2] \le \sigma^2$ and bounded gradients $|g_{t, i}| \le G$. Under a decreasing step-size schedule $\eta_t = \frac{\eta_0}{\sqrt{t}}$, the average squared gradient norm converges as:
\begin{equation}
\min_{t=1, \dots, T} \mathbb{E}\left[\|\nabla \mathcal{L}(\theta_t)\|_2^2\right] \le \mathcal{O}\left(\frac{\ln T}{\sqrt{T}}\right)
\end{equation}
\end{theorem}

\begin{proof}
By $L$-smoothness of $\mathcal{L}$, the descent lemma provides:
\begin{equation}
\mathcal{L}(\theta_{t}) \le \mathcal{L}(\theta_{t-1}) + \langle \nabla \mathcal{L}(\theta_{t-1}), \theta_t - \theta_{t-1} \rangle + \frac{L}{2} \|\theta_t - \theta_{t-1}\|_2^2
\end{equation}
Substituting coordinate update $\Delta \theta_{t, i} = -\frac{\eta_t}{\sqrt{v_{t, i}} + \epsilon} g_{t, i}$:
\begin{equation}
\mathcal{L}(\theta_t) \le \mathcal{L}(\theta_{t-1}) - \sum_{i=1}^P \frac{\eta_t \nabla_i \mathcal{L}(\theta_{t-1}) g_{t, i}}{\sqrt{v_{t, i}} + \epsilon} + \frac{L \eta_t^2}{2} \sum_{i=1}^P \frac{g_{t, i}^2}{(\sqrt{v_{t, i}} + \epsilon)^2}
\end{equation}
Since $|g_{t, i}| \le G$, $v_{t, i} \le G^2$, we bound the denominator $\epsilon \le \sqrt{v_{t, i}} + \epsilon \le G + \epsilon$. Taking conditional expectations:
\begin{equation}
\mathbb{E}[\mathcal{L}(\theta_t) \mid \theta_{t-1}] \le \mathcal{L}(\theta_{t-1}) - \frac{\eta_t}{G + \epsilon} \|\nabla \mathcal{L}(\theta_{t-1})\|_2^2 + \frac{L \eta_t^2 P G^2}{2 \epsilon^2}
\end{equation}
Summing from $t=1$ to $T$ and telescoping yields $\sum_{t=1}^T \eta_t \mathbb{E}[\|\nabla \mathcal{L}(\theta_{t-1})\|_2^2] \le (G+\epsilon)(\mathcal{L}(\theta_0) - \mathcal{L}^*) + \mathcal{O}\left(\sum_{t=1}^T \eta_t^2\right)$. Setting $\eta_t \propto 1/\sqrt{t}$ gives $\sum_{t=1}^T \eta_t = \mathcal{O}(\sqrt{T})$ and $\sum \eta_t^2 = \mathcal{O}(\ln T)$, establishing $\min_{t \le T} \mathbb{E}[\|\nabla \mathcal{L}(\theta_t)\|_2^2] = \mathcal{O}\left(\frac{\ln T}{\sqrt{T}}\right)$.
\end{proof}

\subsection{Limitations and Failure Modes}
\begin{itemize}
    \item \textbf{Non-Monotonicity}: Unlike AdaGrad, $\sqrt{v_t}$ can decrease rapidly when traversing flat plateaus after steep drops, which can cause step sizes to spike and destabilize training unless bounded.
    \item \textbf{Sensitivity to $\alpha$}: If $\alpha$ is too small (e.g., $< 0.8$), variance estimation becomes noisy and fluctuates erratically with single minibatch outliers.
\end{itemize}

\section{Complexity Analysis}

\subsection{Time Complexity}
\begin{itemize}
    \item \textbf{Per-step Update}: For parameter vector of size $P$, computing $v_{\text{next}}$ requires $P$ squaring, multiplication, and addition operations ($\mathcal{O}(P)$).
    \item \textbf{Parameter Shift}: Square root, addition of $\epsilon$, division, and subtraction require $P$ scalar operations ($\mathcal{O}(P)$).
    \item \textbf{Total Time Complexity}: $\Theta(P)$ FLOPs per optimization step.
\end{itemize}

\subsection{Space Complexity}
\begin{itemize}
    \item \textbf{Auxiliary State}: The optimizer maintains exactly one persistent buffer $v \in \mathbb{R}^P$.
    \item \textbf{Auxiliary Memory}: Exactly $P$ floating-point words ($8P$ bytes in FP64), yielding $\Theta(P)$ auxiliary spatial overhead.
\end{itemize}

\section{Worked Numerical Example}

Consider a $2$-dimensional parameter vector $\theta = [1.0, -2.0]^T$, gradient vector $g = [0.2, -0.4]^T$, prior moving average $v = [0.04, 0.09]^T$, learning rate $\eta = 0.01$, $\alpha = 0.9$, and $\epsilon = 10^{-8}$.

\subsection{Coordinate 1 Calculation ($i=1$)}
\begin{enumerate}
    \item $g_1^2 = (0.2)^2 = 0.04$.
    \item $v_{\text{next}, 1} = \alpha v_1 + (1 - \alpha) g_1^2 = 0.9(0.04) + 0.1(0.04) = 0.036 + 0.004 = 0.0400$.
    \item $\sqrt{v_{\text{next}, 1}} + \epsilon = \sqrt{0.04} + 10^{-8} = 0.20000001$.
    \item Step size factor: $\frac{\eta}{\sqrt{v_{\text{next}, 1}} + \epsilon} = \frac{0.01}{0.20000001} \approx 0.0499999975$.
    \item Update: $\theta_{\text{next}, 1} = 1.0 - (0.0499999975)(0.2) = 1.0 - 0.01000000 = 0.990000$.
\end{enumerate}

\subsection{Coordinate 2 Calculation ($i=2$)}
\begin{enumerate}
    \item $g_2^2 = (-0.4)^2 = 0.16$.
    \item $v_{\text{next}, 2} = 0.9(0.09) + 0.1(0.16) = 0.081 + 0.016 = 0.0970$.
    \item $\sqrt{v_{\text{next}, 2}} = \sqrt{0.097} \approx 0.31144823$.
    \item Step size factor: $\frac{0.01}{0.31144823 + 10^{-8}} \approx 0.03210806$.
    \item Update: $\theta_{\text{next}, 2} = -2.0 - (0.03210806)(-0.4) = -2.0 + 0.01284322 = -1.98715678$.
\end{enumerate}

\section{Numerical Considerations and Edge Cases}

\begin{itemize}
    \item \textbf{Epsilon Regularization}: Setting $\epsilon$ inside the square root ($\sqrt{v + \epsilon}$) versus outside ($\sqrt{v} + \epsilon$) yields subtle numerical differences. RMSProp typically uses exterior regularization ($\sqrt{v} + \epsilon$) to prevent division by zero even when $v \to 0$.
    \item \textbf{Denormal Floating Point Numbers}: When gradients are consistently zero for inactive sparse features, repeated multiplication by $\alpha < 1$ causes $v$ to decay exponentially into subnormal (denormalized) FP64/FP32 ranges, inducing CPU microcode slowdowns. Standard implementations flush subnormals to zero.
\end{itemize}

\begin{thebibliography}{9}
\bibitem{tieleman2012} T.~Tieleman and G.~Hinton, ``Lecture 6.5-rmsprop: Divide the gradient by a running average of its recent magnitude,'' \emph{COURSERA: Neural Networks for Machine Learning}, vol.~4, no.~2, pp.~26--31, 2012.
\bibitem{ruder2016} S.~Ruder, ``An overview of gradient descent optimization algorithms,'' \emph{arXiv preprint arXiv:1609.04747}, 2016.
\bibitem{mukkamala2017} M.~C.~Mukkamala and M.~Hein, ``Variants of RMSProp and Adagrad with Logarithmic Regret Bounds,'' in \emph{International Conference on Machine Learning (ICML)}, pp.~2545--2553, 2017.
\end{thebibliography}

\end{document}
